%% Recipes_Detector.tex
%% Created:     Tue Apr  4 16:37:17 2017 by Koehler@I-Mac
%% Last change: 2020-09-04
%%
%% subsection for Detector Recipes
%%
%%%%%%%%%%%%%%%%%%%%%%%%%%%%%%%%%%%%%%%%%%%%%%%%%%%%%%%%%%%%%%%%%%%%%%%%%%%%%
\subsection{Detector calibration recipes}
\label{Sec:detector_calibration}

METIS will have three focal plane detector arrays:
\begin{itemize}
\item One $2\mathrm{k}\times 2\mathrm{k}$ HAWAII2RG detector used for
  LM-band imaging and slit spectroscopy.
\item One $2\mathrm{k}\times 2\mathrm{k}$ GeoSnap (Teledyne) detector
  used for N-band imaging and slit spectroscopy.
\item An array of four $2\mathrm{k}\times 2\mathrm{k}$ HAWAII2RG
  detectors used for LM-band integral-field spectroscopy.
\end{itemize}
This section lists recipes that calibrate detector characteristics
independent of a specific instrument mode. Where \CODE{_det} appears
in FITS keywords of input or product files, it is taken to mean
\CODE{_LM}, \CODE{_N} or \CODE{_IFU} according to the detector
array for which data are being processed.
%------------------------------------------------------------------------------------------------------------------
\subsubsection{\REC*{metis_det_lingain}: Detector linearity and gain determination}
\label{sssec:metis_det_lingain}
\label{rec:metis_det_lingain}
\label{rec:metisdetlingain}

The recipe \REC{metis_det_lingain} determines detector (non-)linearity and absolute detector
gain from a set of flat-field frames taken with the broad-band lamp
over a range of detector exposure times (DITs) and flux levels. The
recipe structure will be similar as for \CODE{detmon_ir_lg} % Not a \REC because it is not our recipe
\cite{detmon-manual}; however, further insight into detector behaviour
(in particular of GeoSnap) may necessitate development of more complex
procedures.

The linearity curve is given by the measured background level as a
function of exposure time for constant illumination. For each pixel
the coefficients of a polynomial fit will be recorded in a
coefficient cube, which can in turn be used to correct for
non-linearity in other recipes. Pixels whose coefficients differ
significantly from the majority of pixels will be marked as bad.

Detector gain is typically computed pixelwise as the slope of a linear
fit of the variance against the mean (or median) values over a set of
frames taken over a range of DITs and illumination levels.  For
mid-infrared detectors that suffer from \ac{ELFN}, e.g.\ the AQUARIUS
detector, this approach does not work.  The GeoSnap is not expected to
show \ac{ELFN}, hence gain determination is probably possible.

The set of calibration frames used for this recipes will include
exposures with WCU window closed (\CODE{LAMP OFF}), which will be used
as `dark' frames that capture thermal emission within the
instrument. This is subtracted from all other exposures in the
sequence.

This satisfies \REQ{METIS-5997}.

\newpage
\begin{recipedef}
  Name:                & \REC{metis_det_lingain} \\
  Purpose:             & Measure detector non-linearity and gain \\
%  Requirements:        & \REQ{METIS-XXXX} \\
%  Type:                & Calibration | Science | Technical \\
%  Templates:           & \TPL{...} \\
  Primary input:       & \RAW{DETLIN_det_RAW} (N$\times$) \\
  Secondary inputs:    & \PROD{BADPIX_MAP_det} (optional) \\
  Matched keywords:    & --- none --- \\
  Parameters:          & \CODE{metis_det_lingain.fitdegree}\newline Degree of polynomial that is fit to linearity correction function, default \texttt{1} \\
                       & \CODE{metis_det_lingain.kappa}\newline Kappa factor for sigma clipping for BPM. Values that deviate more than kappa*sigma from the median linearity are flagged., default \texttt{3} \\
                       & \CODE{metis_det_lingain.linlimit}\newline Limit in ADU above which counts are excluded from both gain calculation and linearity correction function, default \texttt{22000.0} \\
                       & \CODE{metis_det_lingain.truelimit}\newline Limit in ADU to consider as true linear. All flux rates will be corrected to the weighted average flux rate defined by the pixel values below this limit., default \texttt{10000.0} \\
  Algorithm:           & We expect two on and two dark (here named off) images per DIT \\
                       & for the Gain and Linearity calculation, similar to ESO's DETMON. \\
                       & The gain is determined from the slope of the average flux \\
                       & of the on-off maps and the variance of the on-on and off-off maps. \\
                       & A gain correction factor is applied to correct for the inter pixel capacitance. \\
                       & The noise on the gain is derived by resampling with replacement. \\
                       & The linearity is determined following the CRIRES approach where for every pixel the flux rate \\
                       & (ADU/s; including dark current) is calculated and the weighted average of the flux rate \\
                       & at the lowest (most linear) fluxes are taken as the true ADU/s flux rate \\
                       & to which all flux rates should be corrected. \\
                       & This directly defines a correction factor as a function of flux with \\
                       & the most linear fluxes having a correction close to 1. \\
                       & An Nth order polynomial is fit (np.polyfit; taking into account errors) to these correction values as function of flux. \\
                       & We add pixels with linearity coefficients that are significantly different from the mean to the BPM (astropy.stats.sigma\_clip). \\
  Output data:         & \PROD{BADPIX_MAP_det} \\
                       & \STATCALIB{GAIN_MAP_det} \\
                       & \PROD{LINEARITY_det} \\
%  Expected accuracies: & \\
  QC1 parameters:      & \QC{QC GAIN COEFF} \\
                       & \QC{QC GAIN LIN} \\
                       & \QC{QC LIN GAIN MEAN} \\
                       & \QC{QC LIN GAIN RMS} \\
                       & \QC{QC LIN MAX FLUX} \\
                       & \QC{QC LIN MIN FLUX} \\
                       & \QC{QC LIN NUM BADPIX} \\
\end{recipedef}



\begin{figure}[hb]
  \centering
    \def \globalscale {0.700000}
    \fontsize{10}{12}\selectfont
    %% Flowchart of metis_det_lingain, generated by metisp/drld/generate_drld.py -- do not edit

% ADDING NEW DEFINITIONS -------------------------------------------- start
\definecolor{listingbg}{gray}{0.95}
\definecolor{darkgreen}{rgb}{0.0, 0.7, 0.0}
\definecolor{darkblue} {rgb}{0.0, 0.0, 0.7}
\definecolor{cyan} {rgb}{0.0, 0.4, 0.4}
\definecolor{darkred}  {rgb}{0.7, 0.0, 0.0}
\definecolor{darkorange}{rgb}{1.0, 0.49, 0.0}
\definecolor{violet}{rgb}{255, 0, 255}
\definecolor{turq}{rgb}{0.0, 0.7, 0.8}
\definecolor{fits}{rgb}{0.4, 0.1, 1}


\makeatletter
\lstdefinestyle{RAWstyle}{%
  basicstyle=\ttfamily\color{black}%
  \lst@ifdisplaystyle\scriptsize\fi}

\lstdefinestyle{PARstyle}{%
  basicstyle=\ttfamily\color{black}%
  \lst@ifdisplaystyle\scriptsize\fi}

\lstdefinestyle{DRLstyle}{%
  basicstyle=\ttfamily\color{black}%
  \lst@ifdisplaystyle\scriptsize\fi}

\lstdefinestyle{RECstyle}{%
  basicstyle=\ttfamily\color{black}%
  \lst@ifdisplaystyle\scriptsize\fi}

\lstdefinestyle{QCstyle}{%
  basicstyle=\ttfamily\color{black}%
  \lst@ifdisplaystyle\scriptsize\fi}

\lstdefinestyle{TPLstyle}{%
  basicstyle=\ttfamily\color{black}%
  \lst@ifdisplaystyle\scriptsize\fi}

\lstdefinestyle{PRODstyle}{%
  basicstyle=\ttfamily\color{black}%
  \lst@ifdisplaystyle\scriptsize\fi}

\lstdefinestyle{EXTCALIBstyle}{%
  basicstyle=\ttfamily\color{black}%
  \lst@ifdisplaystyle\scriptsize\fi}

\lstdefinestyle{STATCALIBstyle}{%
  basicstyle=\ttfamily\color{black}%
  \lst@ifdisplaystyle\scriptsize\fi}

\lstdefinestyle{FITSstyle}{%
  basicstyle=\ttfamily\color{black}%
  \lst@ifdisplaystyle\scriptsize\fi}
\makeatother

\input{recipe_config}
%% The picture is set in a box and only scaled down when it would not fit the page.
\ifdefined\flowchartbox\else\newsavebox{\flowchartbox}\fi
\sbox{\flowchartbox}{%
\begin{tikzpicture}
  [x=1cm,
  y=-1cm,
  align=center,
  node distance=2cm and 3.5cm,
  redstep/.append style={minimum height=0.8cm, text width=3.8cm}]
  \sffamily
  \pic (start){start};
  \node (input) [below=0.75cm of start-m, input] {%
    N$\times$ \RAW{DETLIN_det_RAW}};

  %% algorithm steps
  \node (step_subtractdark) [below=1.6cm of input, redstep]{%
    subtract dark};
  \node (step_computegain) [below=0.6cm of step_subtractdark, redstep]{%
    compute gain};
  \node (step_linearitycheck) [below=1.25cm of step_computegain, redstep]{%
    linearity check};
  \node (step_thresholding) [below=1.25cm of step_linearitycheck, redstep]{%
    thresholding};
  \pic (stop) [below=1.45cm of step_thresholding] {stop};

  \draw [connection_arrow] (start-m) -- (input);
  \draw [connection_arrow] (input) -- (step_subtractdark);
  \draw [connection_arrow] (step_subtractdark) -- (step_computegain);
  \draw [connection_arrow] (step_computegain) -- (step_linearitycheck);
  \draw [connection_arrow] (step_linearitycheck) -- (step_thresholding);
  \draw [connection_arrow] (step_thresholding) -- (stop-t);

  %% calibration inputs, attached above the step that consumes them
  \node (connect_badpixmap) [connection] at ($(input.south) + (0, 0.97)$) {};
  \node (badpixmap) [left=of connect_badpixmap, calproduct] {\PROD{BADPIX_MAP_det}};
  \draw [connection_arrow, dashed] (badpixmap) -- (connect_badpixmap);

  %% products, attached below the step that produces them
  \node (connect_gainmap) [connection] at ($(step_computegain.south) + (0, 0.62)$) {};
  \node (gainmap) [right=of connect_gainmap, calproduct] {\STATCALIB{GAIN_MAP_det}};
  \draw [connection_arrow] (connect_gainmap) -- (gainmap);
  \node (connect_linearity) [connection] at ($(step_linearitycheck.south) + (0, 0.62)$) {};
  \node (linearity) [right=of connect_linearity, calproduct] {\PROD{LINEARITY_det}};
  \draw [connection_arrow] (connect_linearity) -- (linearity);
  \node (connect_badpixmap) [connection] at ($(step_thresholding.south) + (0, 0.62)$) {};
  \node (badpixmap) [right=of connect_badpixmap, calproduct] {\PROD{BADPIX_MAP_det}};
  \draw [connection_arrow] (connect_badpixmap) -- (badpixmap);

  \draw [frame] ($(input.south) + (-2.85, 0.12)$) rectangle
  ($(step_thresholding.south) + (2.85, 1.0499999999999998)$);
  \node [framecolor, anchor=north west] at ($(input.south) + (-2.85, 0.12)$) {%
    \REC{metis_det_lingain}};
\end{tikzpicture}}%
\ifdim\ht\flowchartbox>\dimexpr\textheight-4\baselineskip\relax
  \resizebox{!}{\dimexpr\textheight-4\baselineskip\relax}{\usebox{\flowchartbox}}%
\else\ifdim\wd\flowchartbox>\linewidth
  \resizebox{\linewidth}{!}{\usebox{\flowchartbox}}%
\else
  \usebox{\flowchartbox}%
\fi\fi

% ADDING NEW DEFINITIONS -------------------------------------------- start
\definecolor{listingbg}{gray}{0.95}
\definecolor{darkgreen}{rgb}{0.0, 0.7, 0.0}
\definecolor{darkblue} {rgb}{0.0, 0.0, 0.7}
\definecolor{cyan} {rgb}{0.0, 0.4, 0.4}
\definecolor{darkred}  {rgb}{0.7, 0.0, 0.0}
\definecolor{darkorange}{rgb}{1.0, 0.49, 0.0}
\definecolor{violet}{rgb}{255, 0, 255}
\definecolor{turq}{rgb}{0.0, 0.7, 0.8}
\definecolor{fits}{rgb}{0.4, 0.1, 1}


\makeatletter
\lstdefinestyle{RAWstyle}{%
  basicstyle=\ttfamily\color{fits}%
  \lst@ifdisplaystyle\scriptsize\fi}

\lstdefinestyle{PARstyle}{%
  basicstyle=\ttfamily\color{cyan}%
  \lst@ifdisplaystyle\scriptsize\fi}

\lstdefinestyle{DRLstyle}{%
  basicstyle=\ttfamily\color{violet}%
  \lst@ifdisplaystyle\scriptsize\fi}

\lstdefinestyle{RECstyle}{%
  basicstyle=\ttfamily\color{darkgreen}%
  \lst@ifdisplaystyle\scriptsize\fi}

%% Write QC parameters like this: \QC{QC_SOMETHING_OR_OTHER}
\lstdefinestyle{QCstyle}{%
  basicstyle=\ttfamily\color{darkblue}%
  \lst@ifdisplaystyle\scriptsize\fi}

%% Write templates like this: \TPL{DARK_LM}
\lstdefinestyle{TPLstyle}{%
  basicstyle=\ttfamily\color{darkred}%
  \lst@ifdisplaystyle\scriptsize\fi}

%% Write products like this: \PROD{SOME_THING}
\lstdefinestyle{PRODstyle}{%
  basicstyle=\ttfamily\color{darkorange}%
  \lst@ifdisplaystyle\scriptsize\fi}

%% external calib files
\lstdefinestyle{EXTCALIBstyle}{%
  basicstyle=\ttfamily\color{Turquoise}%
  \lst@ifdisplaystyle\scriptsize\fi}

% static calib files
\lstdefinestyle{STATCALIBstyle}{%
  basicstyle=\ttfamily\color{teal}%
  \lst@ifdisplaystyle\scriptsize\fi}

% static calib files
\lstdefinestyle{FITSstyle}{%
  basicstyle=\ttfamily\color{black}%
  \lst@ifdisplaystyle\scriptsize\fi}
\makeatother


  \caption[Recipe: \REC*{metis_det_lingain}]{\REC*{metis_det_lingain} --
    determination of linearity and gain of the detectors.}
  \label{Fig:rec_det_lingain}
\end{figure}

%------------------------------------------------------------------------------------------------------------------
\clearpage

\subsubsection{\REC*{metis_det_dark}: Master dark creation}
\label{sssec:metis_det_dark}
\label{rec:det_dark}
\label{rec:metis_det_dark}

Darks are taken in daytime for all science detectors
\cite{METIS-calibration_plan}.  Each dark frame contains a single frame, and calculating a master dark
is done separately for each combination of detector (~\FITS{DET.ID}, \FITS*{DET.CHIP.ID}) and integration time
(\FITS{DET.DIT}, \FITS{DET.NDIT}).\footnote{The dark current is not expected to depend on the readout mode of the detectors. Should hardware tests reveal such a dependence, the recipe will be amended to classify on readout mode as well.} There will be ``METIS-dark''
(with the CLOSED position of the CFO-PP1 wheel) and ``Imager-dark''
(with the CLOSED position in the subsystem PP1), to be distinguished
by header keywords.
% TODO: Use a different DO.CATG for these darks. See https://github.com/AstarVienna/METIS_DRLD/issues/170
The former will be used for pipeline processing, the  latter for monitoring purposes.

For the
IFU, both raw frames and master dark have four extensions
corresponding to the four detectors in the focal-plane array. The
recipe also produces bad pixel masks by identifying hot pixels whose
dark current differs significantly (by more than $\pm 5\sigma$) from
the average over the detector.

This fulfills \REQ{METIS-6063}.

\begin{recipedef}
  Name:                & \REC{metis_det_dark} \\
  Purpose:             & Create master dark \\
%  Requirements:        & \REQ{METIS-XXXX} \\
%  Type:                & Calibration | Science | Technical \\
%  Templates:           & \TPL{...} \\
  Primary input:       & \RAW{DARK_det_RAW} (1--N) \\
  Secondary inputs:    & \PROD{BADPIX_MAP_det} (optional) \\
                       & \STATCALIB{GAIN_MAP_det} (optional) \\
                       & \EXTCALIB{PERSISTENCE_MAP} (optional) \\
  Matched keywords:    & --- none --- \\
  Parameters:          & \CODE{metis_det_dark.stacking.method}\newline Name of the method used to combine the input images (\texttt{average}, \texttt{median}, \texttt{sigclip}), default \texttt{average} \\
                       & \CODE{metis_det_dark.outliers.kappa_low}\newline Lower bound for bad pixel clipping, in standard deviations, default \texttt{5} \\
                       & \CODE{metis_det_dark.outliers.kappa_high}\newline Upper bound for bad pixel clipping, in standard deviations, default \texttt{5} \\
  Algorithm:           & - Group files by detector and \texttt{DIT}, based on header keywords \\
                       & - Call function \texttt{metis\_determine\_dark} for each set of files \\
                       & - Call \texttt{metis\_update\_dark\_mask} to flag deviant pixels \\
  Output data:         & \PROD{MASTER_DARK_det} \\
%  Expected accuracies: & \\
  QC1 parameters:      & \QC{QC DARK MEAN} \\
                       & \QC{QC DARK MEDIAN} \\
                       & \QC{QC DARK MEDIAN MAX} \\
                       & \QC{QC DARK MEDIAN MEAN} \\
                       & \QC{QC DARK MEDIAN MEDIAN} \\
                       & \QC{QC DARK MEDIAN MIN} \\
                       & \QC{QC DARK MEDIAN RMS} \\
                       & \QC{QC DARK NBADPIX} \\
                       & \QC{QC DARK NCOLDPIX} \\
                       & \QC{QC DARK NHOTPIX} \\
                       & \QC{QC DARK RMS} \\
\end{recipedef}



\begin{figure}[hb]
  \centering
        \def \globalscale {0.700000}
        \fontsize{10}{12}\selectfont
        %% Flowchart of metis_det_dark, generated by metisp/drld/generate_drld.py -- do not edit

% ADDING NEW DEFINITIONS -------------------------------------------- start
\definecolor{listingbg}{gray}{0.95}
\definecolor{darkgreen}{rgb}{0.0, 0.7, 0.0}
\definecolor{darkblue} {rgb}{0.0, 0.0, 0.7}
\definecolor{cyan} {rgb}{0.0, 0.4, 0.4}
\definecolor{darkred}  {rgb}{0.7, 0.0, 0.0}
\definecolor{darkorange}{rgb}{1.0, 0.49, 0.0}
\definecolor{violet}{rgb}{255, 0, 255}
\definecolor{turq}{rgb}{0.0, 0.7, 0.8}
\definecolor{fits}{rgb}{0.4, 0.1, 1}


\makeatletter
\lstdefinestyle{RAWstyle}{%
  basicstyle=\ttfamily\color{black}%
  \lst@ifdisplaystyle\scriptsize\fi}

\lstdefinestyle{PARstyle}{%
  basicstyle=\ttfamily\color{black}%
  \lst@ifdisplaystyle\scriptsize\fi}

\lstdefinestyle{DRLstyle}{%
  basicstyle=\ttfamily\color{black}%
  \lst@ifdisplaystyle\scriptsize\fi}

\lstdefinestyle{RECstyle}{%
  basicstyle=\ttfamily\color{black}%
  \lst@ifdisplaystyle\scriptsize\fi}

\lstdefinestyle{QCstyle}{%
  basicstyle=\ttfamily\color{black}%
  \lst@ifdisplaystyle\scriptsize\fi}

\lstdefinestyle{TPLstyle}{%
  basicstyle=\ttfamily\color{black}%
  \lst@ifdisplaystyle\scriptsize\fi}

\lstdefinestyle{PRODstyle}{%
  basicstyle=\ttfamily\color{black}%
  \lst@ifdisplaystyle\scriptsize\fi}

\lstdefinestyle{EXTCALIBstyle}{%
  basicstyle=\ttfamily\color{black}%
  \lst@ifdisplaystyle\scriptsize\fi}

\lstdefinestyle{STATCALIBstyle}{%
  basicstyle=\ttfamily\color{black}%
  \lst@ifdisplaystyle\scriptsize\fi}

\lstdefinestyle{FITSstyle}{%
  basicstyle=\ttfamily\color{black}%
  \lst@ifdisplaystyle\scriptsize\fi}
\makeatother

\input{recipe_config}
%% The picture is set in a box and only scaled down when it would not fit the page.
\ifdefined\flowchartbox\else\newsavebox{\flowchartbox}\fi
\sbox{\flowchartbox}{%
\begin{tikzpicture}
  [x=1cm,
  y=-1cm,
  align=center,
  node distance=2cm and 3.5cm,
  redstep/.append style={minimum height=0.8cm, text width=3.8cm}]
  \sffamily
  \pic (start){start};
  \node (input) [below=0.75cm of start-m, input] {%
    \textsl{N} \RAW{DARK_det_RAW}};

  %% algorithm steps
  \node (step_applynonlinearitycorrect) [below=2.6cm of input, redstep]{%
    apply non-linearity correction};
  \node (step_applypersistencecorrecti) [below=1.25cm of step_applynonlinearitycorrect, redstep]{%
    apply persistence correction};
  \node (step_perpixelmedianmeanfilter) [below=0.6cm of step_applypersistencecorrecti, redstep]{%
    per pixel median/ mean filtering};
  \node (step_thresholding) [below=1.25cm of step_perpixelmedianmeanfilter, redstep]{%
    thresholding};
  \pic (stop) [below=0.9cm of step_thresholding] {stop};

  \draw [connection_arrow] (start-m) -- (input);
  \draw [connection_arrow] (input) -- (step_applynonlinearitycorrect);
  \draw [connection_arrow] (step_applynonlinearitycorrect) -- (step_applypersistencecorrecti);
  \draw [connection_arrow] (step_applypersistencecorrecti) -- (step_perpixelmedianmeanfilter);
  \draw [connection_arrow] (step_perpixelmedianmeanfilter) -- (step_thresholding);
  \draw [connection_arrow] (step_thresholding) -- (stop-t);

  %% calibration inputs, attached above the step that consumes them
  \node (connect_badpixmap) [connection] at ($(input.south) + (0, 0.97)$) {};
  \node (badpixmap) [left=of connect_badpixmap, calproduct] {\PROD{BADPIX_MAP_det}};
  \draw [connection_arrow, dashed] (badpixmap) -- (connect_badpixmap);
  \node (connect_gainmap) [connection] at ($(input.south) + (0, 1.98)$) {};
  \node (gainmap) [left=of connect_gainmap, external] {\STATCALIB{GAIN_MAP_det}};
  \draw [connection_arrow, dashed] (gainmap) -- (connect_gainmap);
  \node (connect_persistencemap) [connection] at ($(step_applynonlinearitycorrect.south) + (0, 0.62)$) {};
  \node (persistencemap) [left=of connect_persistencemap, external] {\EXTCALIB{PERSISTENCE_MAP}};
  \draw [connection_arrow, dashed] (persistencemap) -- (connect_persistencemap);

  %% products, attached below the step that produces them
  \node (connect_masterdark) [connection] at ($(step_perpixelmedianmeanfilter.south) + (0, 0.62)$) {};
  \node (masterdark) [right=of connect_masterdark, calproduct] {\PROD{MASTER_DARK_det}};
  \draw [connection_arrow] (connect_masterdark) -- (masterdark);

  \draw [frame] ($(input.south) + (-2.85, 0.12)$) rectangle
  ($(step_thresholding.south) + (2.85, 0.5)$);
  \node [framecolor, anchor=north west] at ($(input.south) + (-2.85, 0.12)$) {%
    \REC{metis_det_dark}};
\end{tikzpicture}}%
\ifdim\ht\flowchartbox>\dimexpr\textheight-4\baselineskip\relax
  \resizebox{!}{\dimexpr\textheight-4\baselineskip\relax}{\usebox{\flowchartbox}}%
\else\ifdim\wd\flowchartbox>\linewidth
  \resizebox{\linewidth}{!}{\usebox{\flowchartbox}}%
\else
  \usebox{\flowchartbox}%
\fi\fi

% ADDING NEW DEFINITIONS -------------------------------------------- start
\definecolor{listingbg}{gray}{0.95}
\definecolor{darkgreen}{rgb}{0.0, 0.7, 0.0}
\definecolor{darkblue} {rgb}{0.0, 0.0, 0.7}
\definecolor{cyan} {rgb}{0.0, 0.4, 0.4}
\definecolor{darkred}  {rgb}{0.7, 0.0, 0.0}
\definecolor{darkorange}{rgb}{1.0, 0.49, 0.0}
\definecolor{violet}{rgb}{255, 0, 255}
\definecolor{turq}{rgb}{0.0, 0.7, 0.8}
\definecolor{fits}{rgb}{0.4, 0.1, 1}


\makeatletter
\lstdefinestyle{RAWstyle}{%
  basicstyle=\ttfamily\color{fits}%
  \lst@ifdisplaystyle\scriptsize\fi}

\lstdefinestyle{PARstyle}{%
  basicstyle=\ttfamily\color{cyan}%
  \lst@ifdisplaystyle\scriptsize\fi}

\lstdefinestyle{DRLstyle}{%
  basicstyle=\ttfamily\color{violet}%
  \lst@ifdisplaystyle\scriptsize\fi}

\lstdefinestyle{RECstyle}{%
  basicstyle=\ttfamily\color{darkgreen}%
  \lst@ifdisplaystyle\scriptsize\fi}

%% Write QC parameters like this: \QC{QC_SOMETHING_OR_OTHER}
\lstdefinestyle{QCstyle}{%
  basicstyle=\ttfamily\color{darkblue}%
  \lst@ifdisplaystyle\scriptsize\fi}

%% Write templates like this: \TPL{DARK_LM}
\lstdefinestyle{TPLstyle}{%
  basicstyle=\ttfamily\color{darkred}%
  \lst@ifdisplaystyle\scriptsize\fi}

%% Write products like this: \PROD{SOME_THING}
\lstdefinestyle{PRODstyle}{%
  basicstyle=\ttfamily\color{darkorange}%
  \lst@ifdisplaystyle\scriptsize\fi}

%% external calib files
\lstdefinestyle{EXTCALIBstyle}{%
  basicstyle=\ttfamily\color{Turquoise}%
  \lst@ifdisplaystyle\scriptsize\fi}

% static calib files
\lstdefinestyle{STATCALIBstyle}{%
  basicstyle=\ttfamily\color{teal}%
  \lst@ifdisplaystyle\scriptsize\fi}

% static calib files
\lstdefinestyle{FITSstyle}{%
  basicstyle=\ttfamily\color{black}%
  \lst@ifdisplaystyle\scriptsize\fi}
\makeatother


  \caption[Recipe: \REC*{metis_det_dark}]{\REC*{metis_det_dark} -- creation of master
    dark and bad pixel maps}
  \label{Fig:rec_det_dark}
\end{figure}
\clearpage
%------------------------------------------------------------------------------------------------------------------
\subsubsection{\REC*{metis_det_persistence}: Persistence map creation}
\label{sssec:metis_det_persistence}
\label{rec:metis_det_persistence}

Infrared detectors are prone to persistence due to charges trapped on a variety of timescales. The correction for a given science or
calibration exposure is built from a sequence of exposures preceding the exposure in question.
See Section~\ref{ssec:criticalpersistencecorrection} for more details.

As these may include exposures taken for another proprietary programme, the recipe is run by ESO on data taken from the science archive and its products are again ingested into the archive. The recipe will make use of the \ac{HDRL} function for the creation of the persistence map, which will be provided by \ac{ESO} in the future.\\
%We will update this Section as soon as more information is available.

The persistence correction is described in more detail in Section~\ref{ssec:criticalpersistencecorrection}.

%\begin{recipedef}
%  Name:                & \REC{metis_det_persistence}           \\
%  Purpose:             & compute persistence correction maps        \\
%  Requirements:        & \REQ{METIS-9145}                      \\
%  Type:                & Calibration                           \\
%  Templates:           & --                                    \\
%  Parameters:          & \TBD                                  \\
%  Algorithm:           & see hdrl functions:                   \\
%  Output data:         & \PROD{PERSISTENCE_MAP}                \\
%  Expected accuracies: & \TBD                                  \\
%  QC1 parameters:      & see hdrl functions:                   \\
%  hdrl functions:      & \TBD (\CODE{hdrl_persistence_compute} \\
%\end{recipedef}

%\begin{figure}[hb]
%  \centering
%  \resizebox{0.6\textwidth}{0.1\textwidth}{\TODO{\fbox{Figure to be done}}}
%  \caption[Recipe:
%  \REC*{metis_det_persistence}]{\REC*{metis_det_persistence} -- creation
%    of persistence correction frames.}
%  \label{Fig:rec_det_persistence}
%\end{figure}

%%%%%%%%%%%%%%%%%%%%%%%%%%%%%%%%%%%%%%%%%%%%%%%%%%%%%%%



%%% Local Variables:
%%% TeX-master: "METIS_DRLD"
%%% End:
